\subsection{Resolution sensitivity}\label{sec:results-resolution}

The resolution study varies the grid size $K\in\{5,25,49\}$ at $M=10$ and the particle budget $M\in\{5,10,25\}$ at $K=25$ around the primary pair $(25,10)$ at $n=1000$, following the resolution design of Section~\ref{sec:resolution-design}. The pre-fixed retention rule of equation~\eqref{eq:sens-retention-rule} failed for the incumbent pair. On seeds $0,\ldots,9$, the native reference-TCATE error in IC1 is $15.5$ percent smaller at $K=5$ and $14.1$ percent smaller at $K=49$ than at $K=25$, while $K=5$ is $19.1$ percent larger on native law error, and no larger pair improves both targets by more than two between-seed standard errors in both IC1 and IC3. Table~\ref{tab:sensitivity-km} reports the one-factor slices, and the rule returned no automatic selection, so the study reports the complete sensitivity table rather than changing the primary specification silently.

\begin{table}[t]
\centering
\scriptsize
\setlength{\tabcolsep}{3pt}
\caption{Sensitivity of WCF to the quantile grid size $K$ and the particle budget $M$ at $n=1000$ for designs IC0 to IC3. Panel A varies $K$ at $M=10$; Panel B varies $M$ at $K=25$. Entries are mean(SE) over ten replications. The native reference TCATE is evaluated on each cell's own $K$-grid, while the common and interior columns are evaluated on fixed 199-level targets. The primary setting $(K,M)=(25,10)$ appears in both panels.}
\label{tab:sensitivity-km}
\begin{tabular}{llrrrrr}
\toprule
\multicolumn{7}{l}{\textit{Panel A: $K$ sensitivity at $M=10$}} \\
DGP & $K$ & Ref. TCATE native & Law common & Ref. TCATE common & Law interior & Ref. TCATE interior \\
\midrule
IC0 & 5 & 0.0318(0.0027) & 0.2751(0.0030) & 0.0321(0.0030) & 0.0256(0.0007) & 0.0321(0.0031) \\
IC0 & 25 & 0.0286(0.0026) & 0.0669(0.0009) & 0.0286(0.0026) & 0.0178(0.0006) & 0.0292(0.0034) \\
IC0 & 49 & 0.0297(0.0034) & 0.0358(0.0006) & 0.0298(0.0034) & 0.0172(0.0007) & 0.0316(0.0037) \\
IC1 & 5 & 0.0261(0.0026) & 0.2474(0.0028) & 0.0249(0.0027) & 0.0280(0.0008) & 0.0248(0.0029) \\
IC1 & 25 & 0.0309(0.0026) & 0.0605(0.0009) & 0.0290(0.0028) & 0.0195(0.0007) & 0.0317(0.0030) \\
IC1 & 49 & 0.0265(0.0024) & 0.0342(0.0007) & 0.0242(0.0024) & 0.0196(0.0007) & 0.0270(0.0031) \\
IC2 & 5 & 0.0302(0.0036) & 0.2434(0.0029) & 0.0519(0.0051) & 0.0265(0.0006) & 0.0313(0.0036) \\
IC2 & 25 & 0.0310(0.0035) & 0.0605(0.0010) & 0.0342(0.0042) & 0.0198(0.0008) & 0.0315(0.0034) \\
IC2 & 49 & 0.0333(0.0033) & 0.0340(0.0006) & 0.0330(0.0037) & 0.0191(0.0005) & 0.0335(0.0035) \\
IC3 & 5 & 0.0531(0.0059) & 0.2454(0.0027) & 0.0542(0.0058) & 0.0297(0.0011) & 0.0582(0.0058) \\
IC3 & 25 & 0.0474(0.0061) & 0.0621(0.0010) & 0.0505(0.0066) & 0.0228(0.0010) & 0.0506(0.0067) \\
IC3 & 49 & 0.0512(0.0055) & 0.0361(0.0008) & 0.0534(0.0055) & 0.0231(0.0010) & 0.0526(0.0058) \\
\midrule
\multicolumn{7}{l}{\textit{Panel B: $M$ sensitivity at $K=25$}} \\
DGP & $M$ & Ref. TCATE native & Law common & Ref. TCATE common & Law interior & Ref. TCATE interior \\
\midrule
IC0 & 5 & 0.0298(0.0028) & 0.0674(0.0010) & 0.0298(0.0028) & 0.0173(0.0006) & 0.0297(0.0037) \\
IC0 & 10 & 0.0286(0.0026) & 0.0669(0.0009) & 0.0286(0.0026) & 0.0178(0.0006) & 0.0292(0.0034) \\
IC0 & 25 & 0.0280(0.0026) & 0.0661(0.0009) & 0.0282(0.0026) & 0.0174(0.0006) & 0.0295(0.0033) \\
IC1 & 5 & 0.0300(0.0032) & 0.0613(0.0009) & 0.0280(0.0032) & 0.0202(0.0007) & 0.0320(0.0035) \\
IC1 & 10 & 0.0309(0.0026) & 0.0605(0.0009) & 0.0290(0.0028) & 0.0195(0.0007) & 0.0317(0.0030) \\
IC1 & 25 & 0.0281(0.0031) & 0.0602(0.0009) & 0.0259(0.0033) & 0.0197(0.0006) & 0.0289(0.0036) \\
IC2 & 5 & 0.0330(0.0035) & 0.0607(0.0010) & 0.0344(0.0036) & 0.0192(0.0007) & 0.0330(0.0035) \\
IC2 & 10 & 0.0310(0.0035) & 0.0605(0.0010) & 0.0342(0.0042) & 0.0198(0.0008) & 0.0315(0.0034) \\
IC2 & 25 & 0.0339(0.0038) & 0.0608(0.0013) & 0.0345(0.0041) & 0.0198(0.0010) & 0.0331(0.0040) \\
IC3 & 5 & 0.0463(0.0064) & 0.0614(0.0008) & 0.0493(0.0066) & 0.0226(0.0009) & 0.0498(0.0062) \\
IC3 & 10 & 0.0474(0.0061) & 0.0621(0.0010) & 0.0505(0.0066) & 0.0228(0.0010) & 0.0506(0.0067) \\
IC3 & 25 & 0.0487(0.0066) & 0.0618(0.0012) & 0.0526(0.0062) & 0.0233(0.0011) & 0.0534(0.0061) \\
\bottomrule
\end{tabular}
\end{table}


Because native-grid targets change with $K$, the resolution evidence is read together with targets evaluated on the fixed common 199-level and interior level sets of Section~\ref{sec:resolution-design}. The $K=49$ full-range law error reduction of roughly $42$ to $46$ percent is a tail-coverage effect: a $K=25$ fit is continued as a constant outside its fitted support, whereas a $K=49$ fit reaches further into the extreme quantile levels. On the interior levels $0.1$ to $0.9$ the two grids are statistically indistinguishable, with all paired $|t|$ statistics below $1.4$ on both seed samples. The full-range reference-TCATE differences at $K=49$ are $+4.2$ percent in IC0, $-16.4$ percent in IC1, $-3.5$ percent in IC2, and $+5.7$ percent in IC3, and the interior differences are $+8.2$, $-15.0$, $+6.4$, and $+3.9$ percent in the same order, so the reference-TCATE gain is confined to IC1 and is borderline ($t=-2.4$ on the primary seeds and $-2.2$ on the confirmation). The marginal reference effect moves against $K=49$ under poor overlap: on the confirmation sample the IC3 native marginal reference effect is $20.8$ percent worse and the full-range common marginal reference effect is $12.1$ percent worse.

\begin{table}[t]
\centering
\scriptsize
\setlength{\tabcolsep}{3pt}
\caption{Seed-paired resolution differences at $n=1000$ for designs IC0 to IC3, computed as $K=49$ minus $K=25$ at $M=10$. Entries are percentage changes, so a negative value means the larger grid has the smaller error; the value in parentheses is the Monte Carlo standard error of the paired percentage difference in percentage points. The first block uses seeds 0 to 9 from the primary run and the second block uses the fresh-seed confirmation at seeds 20 to 39.}
\label{tab:sensitivity-resolution}
\begin{tabular}{lrrrrr}
\toprule
DGP & Full-range Law & Interior Law & Full-range Ref. TCATE & Interior Ref. TCATE & Ref.-ATE common \\
\midrule
\multicolumn{6}{l}{\textit{Seeds 0 to 9 (primary run)}} \\
\midrule
IC0 & -46.4(1.1) & -3.3(3.2) & 4.2(9.0) & 8.2(8.0) & -26.7(15.2) \\
IC1 & -43.4(1.2) & 0.5(3.2) & -16.4(6.9) & -15.0(5.5) & 22.6(20.5) \\
IC2 & -43.8(0.9) & -3.7(4.0) & -3.5(9.3) & 6.4(9.2) & -37.6(18.7) \\
IC3 & -41.8(1.0) & 1.1(3.0) & 5.7(10.2) & 3.9(8.5) & -6.3(20.1) \\
\multicolumn{6}{l}{\textit{Seeds 20 to 39 (fresh-seed confirmation)}} \\
\midrule
IC0 & -46.5(0.7) & -1.7(2.2) & -1.2(5.3) & -0.8(5.0) & 5.0(10.1) \\
IC1 & -43.7(1.0) & -3.4(2.6) & -7.9(3.6) & -2.6(3.7) & -8.2(10.4) \\
IC2 & -44.1(0.7) & -1.3(2.0) & -10.5(4.9) & -0.6(4.8) & -17.7(8.8) \\
IC3 & -42.8(0.9) & -1.3(2.7) & 0.3(3.7) & 0.8(3.4) & 12.1(6.1) \\
\bottomrule
\end{tabular}
\end{table}


The fresh-seed confirmation (seeds $20$ to $39$, 160 cells, no reselection) reproduced the full-range law gain, with reductions of $46.5$, $43.7$, $44.1$, and $42.8$ percent in IC0 to IC3, and the interior parity, with paired law changes of $-1.7$, $-3.4$, $-1.3$, and $-1.3$ percent. The confirmation rows of Table~\ref{tab:sensitivity-resolution} are a stability check on twenty replications, not an equivalence test.

The factorial extension is allocated only to IC1 and IC3, and its interaction is non-monotone: in IC1 the smallest mean within each row moves from $M=10$ at $K=5$ to $M=25$ at $K=25$ and back to $M=5$ at $K=49$, and none of the larger pairs improves both native targets in both regimes by the two-standard-error margin of the pre-fixed rule. The full grid is reported in Appendix~\ref{app:sensitivity-tables} (Table~\ref{tab:sensitivity-factorial}) and does not alter the primary specification.

The particle-budget panel gives no reason to move away from $M=10$. At $K=25$, increasing to $M=25$ changes the common-grid law error by at most $1.2$ percent, and across the factorial grid by at most $1.5$ percent, while costing four to eight times as much, and decreasing to $M=5$ does not improve the interior targets as a pair and does not deliver a compensating native reference-TCATE gain: the interior reference-TCATE error is worse in IC0, IC1, and IC2 (for example, $0.0320$ against $0.0317$ in IC1), the interior law error is essentially unchanged, and the native reference-TCATE error is worse in IC0 ($0.0298$ against $0.0286$) and statistically indistinguishable in IC1 ($0.0300$ against $0.0309$). The primary pair $(K,M)=(25,10)$ is retained, and $(49,10)$ is reported as a tail-fidelity robustness alternative with the interior parity, the IC3 marginal reference-effect regression, and its $1.6$ times cost disclosed.

\subsection{Assignment mechanisms and propensity models}\label{sec:results-assignment}

Table~\ref{tab:sensitivity-assignment} reports native reference-TCATE errors for the four symmetric assignment mechanisms of Section~\ref{sec:assignment-design}. At $n=500$ the WCF errors are $0.0444$, $0.0452$, $0.0597$, and $0.0461$ under SYM-RANDOM, SYM-LIN, SYM-NL, and SYM-MU, the Causal-DRF errors are $0.0422$, $0.0548$, $0.0889$, and $0.0571$, and the DRF errors are $0.0417$, $0.0428$, $0.0746$, and $0.0501$. At $n=1000$ the WCF errors are $0.0286$, $0.0350$, $0.0417$, and $0.0287$; Causal-DRF records $0.0313$, $0.0392$, $0.0814$, and $0.0461$; and DRF records $0.0275$, $0.0340$, $0.0736$, and $0.0406$. WCF has the smallest native reference-TCATE error under the two nonlinear mechanisms SYM-NL and SYM-MU, is at parity under SYM-LIN and the constant-propensity calibration control SYM-RANDOM, and its margin under the nonlinear mechanisms grows with the sample size.

\begin{table}[t]
\centering
\scriptsize
\setlength{\tabcolsep}{3pt}
\caption{Symmetric and nonlinear assignment designs at $n=500$ and $n=1000$. Entries are native reference-TCATE errors, mean(SE) over ten replications; lower is better. SYM-RANDOM is a constant-propensity calibration control, SYM-LIN uses a linear assignment index, and SYM-NL and SYM-MU use nonlinear assignment with the same symmetric outcome laws.}
\label{tab:sensitivity-assignment}
\begin{tabular}{lrrr}
\toprule
DGP & WCF & Causal-DRF & DRF \\
\midrule
\multicolumn{4}{l}{\textit{Panel A: $n=500$}} \\
\midrule
SYM-RANDOM & 0.0444(0.0040) & 0.0422(0.0045) & 0.0417(0.0043) \\
SYM-LIN & 0.0452(0.0041) & 0.0548(0.0054) & 0.0428(0.0053) \\
SYM-NL & 0.0597(0.0085) & 0.0889(0.0096) & 0.0746(0.0080) \\
SYM-MU & 0.0461(0.0034) & 0.0571(0.0044) & 0.0501(0.0051) \\
\multicolumn{4}{l}{\textit{Panel B: $n=1000$}} \\
\midrule
SYM-RANDOM & 0.0286(0.0028) & 0.0313(0.0032) & 0.0275(0.0037) \\
SYM-LIN & 0.0350(0.0028) & 0.0392(0.0030) & 0.0340(0.0022) \\
SYM-NL & 0.0417(0.0058) & 0.0814(0.0072) & 0.0736(0.0066) \\
SYM-MU & 0.0287(0.0022) & 0.0461(0.0040) & 0.0406(0.0034) \\
\bottomrule
\end{tabular}
\end{table}


The alignment pair isolates the relation between assignment and prognosis rather than the propensity spread: SYM-ALIGN assigns with a function of the prognostic covariate $x_1$, SYM-IRREL assigns with the same function of $x_6$, which enters no outcome surface, and the two propensity variables have the same marginal law by construction. At $n=500$ the paired reference-TCATE difference (aligned minus irrelevant) is $+0.0041$ ($p=0.48$) and the paired law-error difference is $-0.0017$ ($p=0.11$); at $n=1000$ the corresponding differences are $-0.0008$ ($p=0.88$) and $-0.0004$ ($p=0.72$). No significant alignment effect is detected at ten replications.

Table~\ref{tab:sensitivity-propensity} replaces the default logistic propensity with a random-forest factory, the default gradient-boosting factory, and the oracle propensity. The random-forest factory is statistically indistinguishable from the logistic default: its paired reference-TCATE differences range from $-8.1$ to $+1.9$ percent with all $|t|$ below $1.4$ across the two regimes and sample sizes (for example, $0.0383$ against $0.0417$ under SYM-NL at $n=1000$). The default gradient-boosting factory is roughly three to four times worse ($0.1471$ against $0.0417$ under SYM-NL at $n=1000$, and $0.1057$ against $0.0287$ under SYM-MU at $n=1000$). The calibration diagnostic attributes this degradation to the factory rather than to the use of a flexible propensity: the gradient-boosting factory's out-of-fold predictions reach $0.001$ and $0.998$, $2$ to $12$ percent of them are clipped, and its root mean squared error against the true propensity is $0.22$ to $0.23$, compared with $0.05$ to $0.16$ for logistic regression and $0.09$ to $0.10$ for the random forest. The oracle propensity improves on logistic by up to $18.7$ percent under SYM-NL ($0.0339$ against $0.0417$ at $n=1000$) and is a diagnostic benchmark for the efficiency lost to the misspecified logistic model, not a feasible competitor.

\begin{table}[t]
\centering
\scriptsize
\setlength{\tabcolsep}{3pt}
\caption{Propensity-model sensitivity on the SYM-NL and SYM-MU designs at $n=500$ and $n=1000$. Entries are native reference-TCATE errors, mean(SE) over ten replications; lower is better. The columns replace the default logistic propensity with a random-forest factory, a gradient-boosting factory, and the oracle propensity, which is a diagnostic and not a feasible competitor. The outcome learner, grid, and particle budget are held fixed.}
\label{tab:sensitivity-propensity}
\begin{tabular}{lrrrr}
\toprule
DGP & Logistic & Random forest & Gradient boosting & Oracle \\
\midrule
\multicolumn{5}{l}{\textit{Panel A: $n=500$}} \\
\midrule
SYM-NL & 0.0597(0.0085) & 0.0565(0.0063) & 0.1737(0.0187) & 0.0578(0.0049) \\
SYM-MU & 0.0461(0.0034) & 0.0470(0.0033) & 0.1524(0.0165) & 0.0426(0.0042) \\
\multicolumn{5}{l}{\textit{Panel B: $n=1000$}} \\
\midrule
SYM-NL & 0.0417(0.0058) & 0.0383(0.0055) & 0.1471(0.0147) & 0.0339(0.0053) \\
SYM-MU & 0.0287(0.0022) & 0.0291(0.0023) & 0.1057(0.0117) & 0.0261(0.0025) \\
\bottomrule
\end{tabular}
\end{table}


These comparisons assess the stability of the calibrated contrasts to the propensity estimator. They do not establish consistency of the outcome model and do not constitute a double-robustness guarantee at the fixed budgets and ten replications used here.
